\documentclass[preview]{standalone}

\usepackage{amsmath}
\usepackage{amssymb}
\usepackage{stellar}
\usepackage{definitions}

\begin{document}

\id{psicologia-oblio-metamemoria}
\genpage

\section{Oblio, strategie e ricostruzione}

\begin{snippet}{oblio-introduzione}
    La memoria funziona spesso in modo adeguato, ma può essere ostacolata da
    competizione tra ricordi, codifiche poco efficaci o ricostruzioni distorte.
    Lo studio dell'oblio riguarda le condizioni in cui l'acquisizione o il
    recupero delle informazioni diventano difficili.
\end{snippet}

\begin{snippetdefinition}{interferenza-definition}{Interferenza}
    L'\emph{interferenza} si verifica quando ricordi diversi competono tra loro
    e rendono più difficile l'acquisizione o il recupero di informazioni.
\end{snippetdefinition}

\begin{snippetdefinition}{interferenza-proattiva-definition}{Interferenza proattiva}
    L'\emph{interferenza proattiva} si verifica quando informazioni acquisite in
    passato rendono difficile l'acquisizione di nuove informazioni.
\end{snippetdefinition}

\begin{snippetdefinition}{interferenza-retroattiva-definition}{Interferenza retroattiva}
    L'\emph{interferenza retroattiva} si verifica quando l'acquisizione di nuove
    informazioni rende difficile il ricordo di informazioni memorizzate in
    precedenza.
\end{snippetdefinition}

\begin{snippetexample}{cambio-numero-telefono-example}{Cambio di numero di telefono}
    Quando si cambia numero di telefono, all'inizio il vecchio numero può
    rendere difficile memorizzare quello nuovo: è interferenza proattiva. Dopo
    aver ripetuto più volte il numero nuovo, può invece diventare difficile
    ricordare quello vecchio: è interferenza retroattiva.
\end{snippetexample}

\begin{snippetdefinition}{ripasso-elaborativo-definition}{Ripasso elaborativo}
    Il \emph{ripasso elaborativo} è una tecnica di memorizzazione che consiste
    nell'arricchire la codifica di una nuova informazione collegandola ad altro
    materiale significativo.
\end{snippetdefinition}

\begin{snippetexample}{cane-tetto-example}{Cane e tetto}
    Per ricordare la coppia \texttt{Cane-Tetto}, si può costruire un'immagine
    mentale di un cane finito sul tetto per recuperare una pallina. Il legame
    immaginativo rende più facile il recupero della coppia.
\end{snippetexample}

\begin{snippetdefinition}{mnemotecniche-definition}{Mnemotecniche}
    Le \emph{mnemotecniche} sono strumenti che permettono di codificare una
    lunga serie di informazioni associandole a conoscenze familiari già
    codificate in precedenza.
\end{snippetdefinition}

\begin{snippetdefinition}{metodo-loci-definition}{Metodo dei loci}
    Il \emph{metodo dei loci} è una mnemotecnica che permette di ricordare
    l'ordine di una lista di nomi o oggetti associandoli a una sequenza di
    luoghi familiari.
\end{snippetdefinition}

\begin{snippetexample}{metodo-loci-example}{Metodo dei loci}
    Per ricordare una lista della spesa, si può collocare mentalmente ogni
    prodotto lungo il percorso abituale da casa all'università. Il recupero
    avviene poi ripercorrendo mentalmente gli stessi luoghi nello stesso ordine.
\end{snippetexample}

\begin{snippetdefinition}{metamemoria-definition}{Metamemoria}
    La \emph{metamemoria} è la capacità di riflettere sul funzionamento dei
    propri processi mnestici e sulle informazioni che si è certi di possedere.
\end{snippetdefinition}

\begin{snippet}{domande-metamemoria}
    Le domande di metamemoria riguardano la sicurezza di possedere una certa
    informazione, l'attendibilità dei propri ricordi e la valutazione delle
    strategie di recupero utilizzate.
\end{snippet}

\begin{snippetdefinition}{memoria-ricostruttiva-definition}{Memoria ricostruttiva}
    La \emph{memoria ricostruttiva} è un tipo di memoria che ricostruisce le
    informazioni basandosi su forme generali di conoscenza già memorizzate.
\end{snippetdefinition}

\begin{snippetdefinition}{distorsioni-memoria-definition}{Distorsioni della memoria}
    Le \emph{distorsioni della memoria} sono circostanze in cui il ricordo
    ricostruito differisce dall'evento reale.
\end{snippetdefinition}

\begin{snippet}{bartlett-ricostruzione}
    Bartlett individuò tre processi di ricostruzione:
    \begin{itemize}
        \item \emph{livellamento}: semplificazione della storia;
        \item \emph{modellamento}: enfatizzazione di alcuni dettagli;
        \item \emph{assimilazione}: cambiamento dei dettagli.
    \end{itemize}
\end{snippet}

\includesnpt[width=52\%,src=/snippet/static-psychology/reconstructive-memory-example.jpg]{centered-img}

\end{document}
